\section{Targeted Damage: Where Uniform Survival Fails}
\label{sec:targeted}

\subsection{Connectivity under attack}
\label{sec:targeted-structure}

Removing nodes in order of degree or betweenness fragments the graph much faster
than random removal (Figure~\ref{fig:attack}). At $f=0.2$ the giant component
holds 0.797 of nodes under random removal, 0.690 under the degree attack and
0.686 under the betweenness attack. At $f=0.5$ these become 0.493, 0.232 and
0.206. The two attacks behave alike because they target largely the same nodes:
41 of the top 50 nodes are shared, a Jaccard similarity of 0.695. Diseases are
5.7\% of all nodes but 14--16\% of the top 1\% under either ranking
(Appendix~\ref{app:figures}).

\begin{figure}[!htbp]
  \centering
  \includegraphics[width=0.56\linewidth]{act1/primekg_5000_bfs_v1/fig_attack_theory.pdf}
  \caption{Giant component under random node removal (mean of ten orders) and
  under fixed degree and betweenness rankings, with the Callaway prediction for
  the degree attack (dashed).}
  \label{fig:attack}
\end{figure}

The Callaway prediction captures the collapse near $f\approx0.6$. Below that,
however, it lies up to 0.11 above the measured curve. For example, at $f=0.3$ it
gives 0.648 where the graph has 0.546. The theory is accurate to 0.0012 for
stub-matched multigraphs with the same degrees
(Section~\ref{sec:blind-connectivity}). The excess fragility must therefore come
from the requirement that the real graph be simple, or from structure that the
degree sequence does not record.

\subsection{Attributing the excess fragility with null models}
\label{sec:targeted-null}

We rewired the standard graph in three ways. Each keeps every node's degree, so
the same fixed ranking applies to all of them, and each removes a different
part of the remaining structure.
\begin{itemize}[leftmargin=*,itemsep=2pt]
\item \textbf{Full rewiring} \cite{maslov2002specificity}: double-edge swaps that
ignore relation type. Only the degree sequence is kept.
\item \textbf{Typed rewiring}: swaps only within one relation. For drug--protein
and disease--protein edges, the protein ends are exchanged. Every node keeps its
degree in each relation.
\item \textbf{2K rewiring}: swaps only between edge ends whose endpoints have the
same degree. This keeps the joint degree distribution, and hence all
degree--degree correlations.
\end{itemize}
Each model was realized five times, with ten attempted swaps per edge. Every
rewired graph is simple and has the original degree sequence. The rewiring is
thorough. After full, typed and 2K rewiring, only 9\%, 13\% and 14\% of the
original edges remain, and the top 1\% of hubs keep 20\%, 27\% and 29\% of their
original neighbors. Global clustering (transitivity) is 0.069 in the real graph
and 0.051--0.070 in the null models. The standard deviation of $S$ across realizations
is at most 0.007 for 2K rewiring and 0.015 for typed rewiring, and reaches 0.046
for full rewiring only near the collapse at $f\approx0.55$.

\begin{figure}[!htbp]
  \centering
  \includegraphics[width=\linewidth]{act1/primekg_5000_bfs_v1/fig_attack_null_model.pdf}
  \caption{Degree attack on the standard graph and on three degree-preserving
  null models. (a) Giant component. (b) Difference from the Callaway prediction.
  Null models are means over five realizations. Shading shows one standard
  deviation.}
  \label{fig:null}
\end{figure}

For $f\le0.5$, Figure~\ref{fig:null} orders the models as expected. Full
rewiring is closest to the theory, typed rewiring is lower, 2K rewiring is lower
still, and the real graph is lowest. We split the gap between theory and the real graph into three
steps. The step from theory to full rewiring is the cost of requiring a simple
graph with a heavy-tailed degree sequence. The step from full to 2K rewiring is
the effect of degree correlations beyond those that the simple-graph constraint
already forces (fully rewired graphs have $r=-0.108$). The step from 2K rewiring to the real graph
is everything else. For $0.1\le f\le0.4$, these steps account for 6--18\%,
32--49\% and 44--53\% of the gap. As $f$ grows, the share of degree correlations
falls and the share of the simple-graph constraint rises.

The mechanism behind the degree-correlation step is visible directly. In the
sample, low-degree nodes hang from hubs (Table~\ref{tab:graphs}), so removing
the hubs strands them. A 5\% attack isolates 6.7\% of all nodes in the real
graph, 4.3\% in the 2K graphs and 1.1\% after full rewiring. Drugs are the
extreme case. After the attack, 77\% of the surviving drugs have no neighbor
left in the real graph and 78\% in the 2K graphs, but only 26\% after full
rewiring. The real graph is more clustered than
the 2K graphs (transitivity 0.069 against 0.051). Clustering is therefore one
candidate for the residual half, together with mesoscale structure such as
disease modules and the BFS sampling itself. We did not separate these.

To separate PrimeKG from its sample, we repeated the degree attack on the full
filtered graph. It also falls below its Callaway prediction, by up to 0.06 near
$f=0.25$ (0.352 against 0.415), before both collapse near $f\approx0.35$. At equal
$f$, the sample's absolute gap is only 1.3--1.4 times larger (0.081 against
0.062 at $f=0.2$). Relative to the predicted value, the full graph is as fragile
or more so (a 12\% shortfall against 10.5\% at $f=0.2$). The excess fragility
under attack is therefore a property of PrimeKG, not an artefact of the sample,
even though the two graphs differ in their degree correlations.

\subsection{Retrieval under attack}
\label{sec:targeted-retrieval}

The node formula, Eq.~\eqref{eq:r2-node}, predicts that 87.9\% of queries
survive the blind removal of 5\% of nodes. After the
degree attack, 30.5\% survive, and after the betweenness attack 24.0\%. The
attack penalties are therefore 57.4 and 63.9 percentage points relative to the
formula, or 56.5 and 63.0 points relative to measured random removal. Random removal of
the same number of nodes leaves 87.0\% (ten-order mean), consistent with the
formula within sampling error.

Equation~\eqref{eq:r2-node} is a product of an endpoint factor $\phi^2$ and a
route factor. The penalty therefore splits into two channels
(Figure~\ref{fig:channels}). A sequential split depends on which factor is
changed first, so we average both orders. For the 5\% degree attack, lost
endpoints account for 17.4 points and lost routes for 40.0. Depending on the
order, these ranges are 11--24 and 34--46.
\begin{itemize}[leftmargin=*,itemsep=2pt]
\item \textbf{Endpoints.} Hub diseases are removed early. The disease of 27.5\%
of queries is in the top 5\% by degree, against 7.5\% for drugs.
\item \textbf{Routes.} A shared protein must link both a drug and a disease, so
the benchmark's routes run through high-degree proteins. Their median degree
is 188 per route, against 25 for all proteins. Of the 462 routes, 61\% pass
through a protein in the top 5\% by degree (degree at least 153). The 5\%
attack removes the only shared protein of 84\% of single-path queries. It also fails 53\% of multi-path
queries: 36\% lose an endpoint, and 17\% keep both endpoints but lose every
route.
\end{itemize}

\begin{figure}[!htbp]
  \centering
  \includegraphics[width=0.62\linewidth]{act2/primekg_5000_bfs_v1/fig_corr_channels.pdf}
  \caption{Attack penalty, the blind-damage prediction minus measured two-hop
  success, split into lost endpoints and lost routes. Each split averages the
  two orders of decomposition.}
  \label{fig:channels}
\end{figure}

Path count is therefore not effective redundancy. The routes of a query pass
through different proteins, so they remain independent. Under an attack,
however, a route survives only if its protein's degree is below the cut-off, and
most routes fail that test. Equation~\eqref{eq:r2-attack}, which replaces the
uniform survival probability by the degree-dependent $\sigma_k$, reproduces the
measured success exactly at 5\% and 10\% damage (0.305 and 0.170). At 20\%, where
ties at the cut-off matter, it gives 0.034 against a measured 0.035. Nothing
beyond the degrees of the endpoints and the route proteins is needed to explain
the penalty. The structural result of Section~\ref{sec:targeted-null} is
analogous: the units that protect a low-degree node, its edges, lead to hubs.
For the same reason, degree features predict which queries fail under attack. A
logistic regression on degrees reaches an AUC of 0.88 in cross-validation over
queries, against 0.66 for the degree-blind formula.

\subsection{What the giant component misses}
\label{sec:targeted-proxy}

At equal damage, the giant component registers that attacks are worse, but it
understates by how much (Figure~\ref{fig:summary}b). After 5\% node removal, the
degree attack lowers $S$ by 6.6 points more than random removal, and lowers
$R_2$ by 56.5 points more. The two metrics also rank the attacks differently. By
$S$ the degree attack is marginally worse (0.883 against 0.886). By $R_2$ the
betweenness attack is worse by 6.5 points (24.0\% against 30.5\%). Beyond 10\%
damage the two attacks differ by at most two of the 200 queries. For a retrieval
task, the useful diagnostic is the redundancy of the task's own routes and the
degrees of the nodes they pass through, not the size of the giant component.
